\documentclass[a4paper,14pt,fleqn]{scrartcl}%fleqn pour aligner les equations à gauche
\usepackage{../../Styles/Cours}


\newcommand{\chnum}{Chapitre 11}
\newcommand{\chtitle}{Mobile autoporteur}
\newcommand{\dtype}{TP}

\begin{document}
\maketitle
Un mobile autoporteur est un dispositif de laboratoire permettant de simuler le déplacement d'un objet qui ne serait pas soumis à des forces de frottements. Il utilise d'une part un coussin d'air qui lui permet de léviter audessus d'une surface plane et d'autre part un système de marquage lui permettant de tracer un point à intervalles de temps réguliers.
\section{Mouvement rectiligne}
\begin{enumerate}
    \item Activer le mobile autoporteur et produire une trajectoire rectiligne.
    \item Sur le tracer, nommer une série de points de 10 points consécutifs $M_1$, ..., $M_{10}$
    \item Tracer les vecteurs déplacements $\vv{M_3M_6}$ et $\vv{M_7M_{10}}$ en vert.
    \item Calculer les valeurs de vitesse $v_5$ et $v_9$
    \item Tracer les vecteurs vitesse $\vv{v_5}$ et $\vv{v_9}$ en rouge. (préciser l'échelle utilisée).
    \item Décrire le mouvement.
\end{enumerate}

\section{Mouvement circulaire}
\begin{enumerate}
    \item Proposer un schéma permettant de produire une trajectoire circulaire avec le mobile autoporteur. Puis réaliser l'expérience.
    \item Sur le tracer, nommer une série de points de 10 points consécutifs $M_1$, ..., $M_{10}$
    \item Tracer les vecteurs déplacements $\vv{M_3M_6}$ et $\vv{M_7M_{10}}$ en vert.
    \item Calculer les valeurs de vitesse $v_5$ et $v_9$
    \item Tracer les vecteurs vitesse $\vv{v_5}$ et $\vv{v_9}$ en rouge. (préciser l'échelle utilisée).
    \item Décrire le mouvement.
\end{enumerate}

\end{document}